\subsubsection{Robustness to Differential Misclassification}
\label{app:differential_stress}

Our theoretical framework assumes non-differential misclassification: $P(A^*{\mid}A, Y) = P(A^*{\mid}A)$. To assess sensitivity to violations of this assumption, we conduct a stress test on the exposure topology (the most vulnerable to sign reversal).

We introduce outcome-dependent differential error by splitting the population at the median of $Y$: units with $Y$ above the median are misclassified according to $C_{\text{high}} = \text{normalize}(C + \varepsilon \cdot \Delta)$, and units below receive $C_{\text{low}} = \text{normalize}(C - \varepsilon \cdot \Delta)$, where $\Delta \sim \mathcal{N}(0,1)^{K \times K}$ is a random perturbation matrix and normalization projects each column onto the probability simplex. The parameter $\varepsilon$ controls the degree of differential error: $\varepsilon{=}0$ recovers the non-differential case.

For each $\varepsilon \in \{0, 0.01, 0.05, 0.10\}$, we sample 10{,}000 random $K{=}3$ confusion matrices from a Dirichlet$(1,1,1)$ prior and simulate $N{=}10{,}000$ observations per matrix. Table~\ref{tab:differential_stress} reports the results.

\begin{table}[h]
\centering
\caption{Exposure topology sign-flip and violation rates under increasing differential misclassification ($\varepsilon$). Sign-flip rate is stable across all $\varepsilon$ levels; amplification increases modestly.}
\label{tab:differential_stress}
\small
\begin{tabular}{cccc}
\toprule
$\varepsilon$ & \textbf{Sign Flip} & \textbf{Amplification} & \textbf{Violation} \\
\midrule
0.00 (non-diff.) & 83.0\% & 13.9\% & 84.7\% \\
0.01 & 83.4\% & 14.0\% & 85.0\% \\
0.05 & 83.7\% & 15.3\% & 85.8\% \\
0.10 & 83.2\% & 18.6\% & 86.4\% \\
\bottomrule
\end{tabular}
\end{table}

The sign-flip rate remains effectively constant (83.0--83.7\%) across all tested levels of differential error, indicating that sign reversal in the exposure topology is driven primarily by the confusion matrix structure $C$ rather than by the differential mechanism. Amplification increases modestly from 13.9\% to 18.6\% at $\varepsilon{=}0.10$, reflecting the additional bias introduced by outcome-dependent misclassification. These results suggest that our non-differential framework provides a conservative characterization of sign-flip risk: relaxing the non-differential assumption does not materially increase the probability of qualitative inferential errors.
